% ID: MAT-GAN-ELL-001
% Titolo: Equazione di un'ellisse dati vertice e fuoco
% Disciplina: Matematica
% Area: Geometria analitica
% Argomento: Ellisse
% Sottoargomento: Ellisse con assi paralleli agli assi cartesiani
% Classe: Quarta liceo scientifico
% Difficolta: 3
% Tipo: problema_numerico
% Risultato: \frac{x^2}{25}+\frac{y^2}{16}=1
% Tempo_stimato: 10 min
% Competenze: rappresentazione analitica, uso dei parametri, calcolo algebrico
% Prerequisiti: definizione di ellisse, semiassi, fuochi
% Tag: matematica, geometria_analitica, ellisse, coniche, fuochi
% Fonte: esempio_originale
% Autore: Riccardo Carli
% Licenza: CC-BY-SA-4.0

\begin{esercizio}
Determina l'equazione dell'ellisse con centro nell'origine, asse maggiore
sull'asse $x$, un vertice in $A(5,0)$ e un fuoco in $F(3,0)$.
\end{esercizio}

\begin{soluzione}
Per un'ellisse con centro nell'origine e asse maggiore sull'asse $x$,
l'equazione ha forma:
\[
\frac{x^2}{a^2}+\frac{y^2}{b^2}=1,
\qquad a>b>0.
\]

Il vertice $A(5,0)$ indica che $a=5$, quindi $a^2=25$.
Il fuoco $F(3,0)$ indica che $c=3$, quindi $c^2=9$.

Per l'ellisse vale:
\[
c^2=a^2-b^2.
\]

Quindi:
\[
b^2=a^2-c^2=25-9=16.
\]

L'equazione cercata e:
\[
\frac{x^2}{25}+\frac{y^2}{16}=1.
\]
\end{soluzione}

